% ============================================================
%  Photogrammetry Model Creation | METIL Sim Jam
%  Authors: Gavin Heaver, Sebastian Noel, Stevin George
%
%  Palette: black body, navy for section headings, orange for
%  placeholders / TODO blocks / figure stubs.
% ============================================================
\documentclass[11pt]{article}

% ---------- Packages ----------
\usepackage[utf8]{inputenc}
\usepackage[T1]{fontenc}
\usepackage[margin=1in]{geometry}
\usepackage[protrusion=true,expansion=false]{microtype}
\usepackage{graphicx}
\usepackage{xcolor}
\usepackage{enumitem}
\usepackage{titlesec}
\usepackage[most]{tcolorbox}
\usepackage{hyperref}
\usepackage{parskip}

% ---------- Colors (black, navy, orange only) ----------
\definecolor{navy}{HTML}{14213D}      % large headings
\definecolor{todoclr}{HTML}{DD6B20}   % orange: placeholders / TODO
\definecolor{todobg}{HTML}{FCEFE3}
\definecolor{iobg}{HTML}{F2F2F2}      % neutral gray fill for I/O boxes

\hypersetup{
    colorlinks=true,
    linkcolor=black,
    urlcolor=black,
    citecolor=black,
    pdftitle={Photogrammetry Model Creation - METIL Sim Jam},
    pdfauthor={Gavin Heaver, Sebastian Noel, Stevin George}
}

% ---------- Heading styling ----------
% Section (large heading) = navy; subsection / subsubsection = black.
\titleformat{\section}
  {\normalfont\Large\bfseries\color{navy}}{\thesection}{0.6em}{}
\titleformat{\subsection}
  {\normalfont\large\bfseries}{\thesubsection}{0.6em}{}
\titleformat{\subsubsection}
  {\normalfont\normalsize\bfseries}{\thesubsubsection}{0.6em}{}

% ---------- Input / Output stage box (neutral) ----------
\newtcolorbox{stageio}{
    enhanced, breakable,
    colback=iobg, colframe=black,
    boxrule=0.4pt, arc=2pt,
    left=8pt, right=8pt, top=6pt, bottom=6pt,
    fontupper=\small
}
\newcommand{\inout}[2]{%
    \begin{stageio}%
    \textbf{Input:}~#1\\[2pt]%
    \textbf{Output:}~#2%
    \end{stageio}%
}

% ---------- TODO / placeholder block (orange) ----------
\newtcolorbox{todobox}[1]{
    enhanced, breakable,
    colback=todobg, colframe=todoclr,
    boxrule=0.6pt, arc=2pt,
    left=8pt, right=8pt, top=6pt, bottom=6pt,
    fonttitle=\bfseries\small,
    coltitle=white, colbacktitle=todoclr,
    title={#1}
}

% ---------- Figure placeholder (orange) ----------
% Renders a labeled box so the document compiles without images.
% Replace \figph{...} with \includegraphics later.
\newcounter{figph}
\newcommand{\figph}[2][0.62\linewidth]{%
  \begin{center}
  \setlength{\fboxsep}{0pt}%
  \fcolorbox{todoclr}{todobg}{%
    \begin{minipage}[c][2.6cm][c]{#1}
      \centering\itshape\color{todoclr}\small [\,Figure placeholder: add image here\,]
    \end{minipage}}%
  \\[3pt]
  \refstepcounter{figph}%
  {\small\itshape\color{todoclr}Figure~\thefigph. #2}
  \end{center}
}

% ---------- Pipeline overview stage (orange placeholder, black label + caption) ----------
\newcommand{\pipestage}[2]{%
  \begin{minipage}[c]{0.28\linewidth}
    \centering
    \setlength{\fboxsep}{0pt}%
    \fcolorbox{todoclr}{todobg}{%
      \begin{minipage}[c][2.1cm][c]{\linewidth}
        \centering\itshape\color{todoclr}\footnotesize [\,Image placeholder\,]
      \end{minipage}}\\[5pt]
    \textbf{#1}\\[1pt]
    {\footnotesize\itshape #2}
  \end{minipage}%
}

\setcounter{tocdepth}{2}   % TOC lists sections and subsections

% ============================================================
\begin{document}

% ---------- Page 1: title, authors, abstract ----------
\begin{center}
    {\LARGE\bfseries\color{navy} Photogrammetry Model Creation}\\[4pt]
    {\large METIL \textbar{} Sim Jam}\\[10pt]
    {\large Gavin Heaver \quad\textbullet\quad Sebastian Noel \quad\textbullet\quad Stevin George}\\[6pt]
\end{center}

\vspace{2pt}
\hrule height 0.8pt
\vspace{14pt}

\section*{Abstract}
\begin{todobox}{Draft, review and refine}
The following is a starting draft assembled from the project notes. Tighten the wording and adjust the claims once the full pipeline (RealityScan, rigging, and Unity import) is documented.
\end{todobox}

\noindent This project documents a repeatable pipeline for turning ordinary, flat (2D) photographs into textured, colorized 3D models that are ready for use in VR and simulation applications. Using RealityScan as the core photogrammetry tool, the process moves from photography, to model generation, to rigging in Blender, and finally to import into Unity. A deliberate constraint of the project was to avoid AI- and cloud-based tools, which act as security risks, while still keeping the workflow convenient and the resulting models lightweight enough to remain computationally inexpensive. We found that high-quality results are achievable with accessible equipment (a single DSLR/mirrorless camera, controlled lighting, and disciplined photography), provided a consistent set of rules is followed. The workflow is organized by object scale (small/medium objects, large objects, extra-large rooms, and people) so that it can serve as a practical, step-by-step guide from the first photo to a Unity-ready model.

\vspace{16pt}
\begin{center}
{\small\bfseries Pipeline Overview}\\[6pt]
\pipestage{Photography}{Capture and organize photos}%
\hfill\raisebox{0.95cm}{\large$\longrightarrow$}\hfill
\pipestage{RealityScan}{Generate the initial 3D model}%
\hfill\raisebox{0.95cm}{\large$\longrightarrow$}\hfill
\pipestage{Blender}{Clean, decimate, and rig}%
\\[6pt]
{\footnotesize\itshape The result feeds into Unity for VR/simulation use.}
\end{center}

\newpage

% ---------- Page 2: table of contents ----------
\tableofcontents

\newpage

% ============================================================
\section{Introduction (Precursors)}
This section covers the goals, limitations, and general principles that apply to the entire pipeline. Read it before attempting any scan; most photography mistakes trace back to one of the principles below.

\subsection{Goals and Major Limitations}
\begin{itemize}[leftmargin=1.4em]
  \item \textbf{Major limits:} Avoid using AI and avoid using the cloud, as both act as security risks.
  \item \textbf{Major goals:} Maintain a good level of convenience (not requiring insanely professional materials) without sacrificing quality, while also keeping the renderings from being computationally expensive (i.e., not an insane number of polygons).
  \item \textbf{Items for photography:} one DSLR camera (with reasonable accessories such as a lens, SD card, battery, etc.), a tripod, a table to support items, lights, and computers for moving images.
\end{itemize}

\subsection{Notes on RealityScan}
\begin{itemize}[leftmargin=1.4em]
  \item RealityScan does support LiDAR files to develop renderings; however, the goal of this project is to limit it to flat, 2D photos to make 3D models. Using LiDAR, either by itself or alongside 2D photos, is a direction for further research (see Section~\ref{sec:future}).
  \item There is a RealityScan mobile app, but we can't use it primarily because its processing comes from Epic Games' cloud servers.
  \item RealityScan is a very smart program. While there are many rules you can follow, there are also many that can be bent.
\end{itemize}

\subsection{General Limitations}
\begin{itemize}[leftmargin=1.4em]
  \item \textbf{Holes:} It doesn't like holes (screw holes or ports). More precisely, it won't scan that depth; it just makes the port look flat.
  \item \textbf{Shine:} It really does not like any type of shine on objects. Truly try to make objects matte.
  \item \textbf{Mirrors / reflective materials:} RealityScan looks for consistency across photos, where objects keep the same appearance from different angles, because that consistency is how it connects points. Reflective materials look different depending on the viewing angle, which ruins that consistency and makes photos that would normally line up no longer share the same 2D details.
\end{itemize}

\subsection{Photography Principles for Good Alignment}
\begin{itemize}[leftmargin=1.4em]
  \item \textbf{Overlap and coverage:} The more photos that follow the rules, the better. A good rule of thumb is 60--80\% overlap between photos (how much of one picture is repeated in another). Another is that any single point on an object should appear in at least 5 different photos.
  \item \textbf{Transitioning between layers:} The same overlap logic applies when adding layers (height/angle bands) of photos. Don't increase or decrease height too much, or the overlap drops and RealityScan sees too many new points it can't place. For example, if the highest angled layer is 30\textdegree{} above the object, add a 15\textdegree{} layer to transition between the middle and the highest, connecting them. Otherwise, when aligning, RealityScan might not use all the photo layers because it can't find points similar enough to attach a layer, effectively wasting that layer, with its detail left out of the model.
  \item \textbf{Each photo should add new data:} ``The more photos, the better'' holds only if each photo shows the object in at least a slightly new way. If two photos are nearly identical, one will be useless or they may slightly confuse the program for that angle.
  \item \textbf{Consistency:} Don't continue if things change for the object being photographed; RealityScan likes consistency. If things don't line up perfectly, it doesn't truly know where they go. If something changes, restart (e.g., while shooting a rock, if it wasn't perfectly balanced and it moves, restart from the new position).
\end{itemize}

\subsection*{Final Note}
Everything here is based on our experience and research. While we tried to cover as much as possible, there are still things we couldn't look into due to limitations; those are covered in the Future Research section (Section~\ref{sec:future}).

% ============================================================
\section{Common Workflow (Shared Across All Categories)}
\label{sec:common}
The steps below are identical for every object-size category. Each category section that follows references this one and lists \emph{only} what differs. When a category section is silent on a topic, the common workflow here applies.

% ----------------------------------------------------------
\subsection{Photography, Common Setup}
\label{sec:common-photo}
\inout{a subject, camera, and lighting to be set up}{organized photos of the subject, exported as JPEG and ready for RealityScan}

\subsubsection{Camera}
Use a DSLR/mirrorless camera if available for the best results. A phone camera is convenient but does not produce results that are as good.

\subsubsection{Lighting}
Use flat, even lighting surrounding the subject. A good arrangement is a light above plus 2--3 lights surrounding the subject (one in front, one in back, and one off to the side to show depth). Take photos inside rather than outside if possible. If outside, shoot on overcast days and keep the light as controlled as possible. The overall goal is to show depth while minimizing shadows and unneeded highlights, so as to capture as much of the subject as possible.

\subsubsection{Camera Settings}
\begin{description}[leftmargin=1.6em, style=nextline]
  \item[Flash] Avoid flash and rely on the flat lighting. Only use it if limited lighting is available and you need it for depth.
  \item[Stability] A tripod gives the absolute best results where available; however, stable hands work perfectly well. Every scan in this project came from stable hands, patience, and confirming that the subject was framed correctly.
  \item[Camera mode] As long as intense filters are avoided and flash isn't used, many standard camera modes are fine. A good option is the equivalent of Canon's ``P Mode'' (Program AE / Smart Auto), which avoids flash while automatically setting shutter speed and aperture; its biggest weakness is glare on shiny objects. For the best results and full control, Manual is recommended. A good baseline is: aperture between f/8 and f/11, ISO 100, shutter speed 1/125\,s, and a manual white balance.
  \item[Low-light compromises] If the light isn't great or details look too dark, adjust in this order to add brightness while minimizing warping or blur: (1) lower the shutter speed (fine on a tripod, but not recommended below 1/40--1/60\,s handheld, or photos blur); (2) raise the ISO; (3) if still needed, widen the aperture to around f/6. Play with these to find the best combination for the situation, starting with shutter speed, then ISO, and lastly aperture.
\end{description}

\noindent\textit{Example (exposure compromise).} The Rubik's cube is a relatively shiny object shot handheld without the best lighting. It was coated with dry shampoo and shot at shutter speed 1/40\,s, aperture f/8, and ISO 1600; the ISO is roughly at the upper limit of what we'd recommend.

\subsubsection{Focus}
Keep the subject in focus in every photo, and try to keep it equally focused across the frame so RealityScan can identify the characteristics surrounding each ``center'' point. A little focus falloff on one point won't ruin a scan, but limit blurriness on the subject as much as possible.

\subsubsection{Filling the Frame}
The frame should completely capture the subject, with the subject filling roughly 70--80\% of the frame. Position yourself based on your focal length to best meet this. Photos can be horizontal or vertical, whichever maximizes the subject in the frame.

\subsubsection{Reducing Shine and Glare}
\label{sec:common-shine}
Coat flat, shiny, transparent, or single-colored objects with scanning spray or dry shampoo to give a matte finish for RealityScan to ``bite onto'' and notice details better. We used dry shampoo because it was what was available, but KCNSC has access to approved scanning spray, so try that first. If the shine itself is important, you can still photograph it, but expect extra time avoiding the shine reading as specific colors or light (e.g., reflecting a surface color); make sure surroundings don't heavily conflict (lights) while also not blending in (e.g., a gray table behind a shiny black object). If the surface is more ``mirror'' than just shiny, give it a matte finish, or scans won't appear correctly (they look concave). If color doesn't matter and you just need a model, you can also add differently colored marks with paint, markers, or tape. A more professional option, researched but not tested here, is using polarizing film over the lights and/or a polarizing filter on the lens; orienting the two 90\textdegree{} apart is called ``cross-polarization'' and eliminates glare. This is discussed further in Section~\ref{sec:future}.

\subsubsection{File Handling}
\begin{itemize}[leftmargin=1.4em]
  \item \textbf{Shoot RAW, but don't use the RAW files.} Shoot in RAW to maximize data, but don't use the RAW photo files in RealityScan.
  \item \textbf{Use high-quality JPEG.} RealityScan supports many file types, but the most consistent results came from JPEG. DSLR/mirrorless RAW files were inconsistent (to the point of crashing the program), so aim for high-quality JPEG.
  \item \textbf{Convert phone photos to JPEG (not HEIC/HEIF).} On iPhone, the simplest direct conversion is to copy all the photos in the Gallery app and paste them into a designated folder in the Files app; this copy-and-paste converts them to JPEG, after which they can be saved to a USB flash drive or emailed. On Android, find a guide for your specific phone, but as a rule of thumb you can either turn off HEIF in the Camera settings (under ``Advanced''), or convert existing photos in the Gallery app by tapping the photo, hitting edit (a pencil icon, or via the three-dots menu), and saving a copy as JPEG.
  \item \textbf{Divide top and bottom files.} Where applicable, split the files into a ``Top'' folder and a ``Bottom'' folder for each object so RealityScan can render them separately before combining. Otherwise it can't orient what's on top versus bottom relative to the surface the object is laying on.
\end{itemize}

% ----------------------------------------------------------
\subsection{RealityScan, Common Workflow}
\label{sec:common-rs}
\inout{organized photos of the subject}{an initial 3D model (textured and colorized)}
\begin{todobox}{To be documented}
Document the shared RealityScan workflow that applies to every category here: importing/aligning photos, combining the top and bottom components, generating the model, texturing/colorizing, and exporting. Category-specific differences go in each category's own RealityScan subsection.
\end{todobox}

% ----------------------------------------------------------
\subsection{Rigging in Blender, Common Workflow}
\label{sec:common-rig}
\inout{an initial 3D model (textured and colorized)}{a cleaned, decimated, and rigged model ready for Unity import}
\begin{todobox}{To be documented}
Document the shared Blender workflow here: cleanup, decimation (polygon reduction toward the low-poly goal), texture/material handling, and rigging. Category-specific differences go in each category's own Rigging subsection.
\end{todobox}

% ============================================================
\section{Small / Medium Objects}
\textit{Examples: a pack of gum, a shoe box.} Follow the Common Workflow (Section~\ref{sec:common}); only the differences are listed below.

\subsection{Photography}
\inout{camera and lighting set up per Section~\ref{sec:common-photo}}{organized photos of the object, split into Top and Bottom folders}

\textbf{Focal length:} 35--50\,mm, with around 35\,mm used most often and giving good results. There's no need to zoom much beyond this.

\paragraph{Capturing the object.} In both methods below, keep the object locked in place, fully inside the frame (no part cut off), and easy to flip. For example, a rock is always stable and can be flipped into another stable position, whereas a paper object can be shifted by a gust of wind, causing inconsistencies in the model.

\textbf{Method 1 (recommended).}
\begin{enumerate}[leftmargin=1.6em]
  \item Walk around the object with the camera level to it, taking a photo about every 5--10\textdegree{} around the full 360\textdegree.
  \item Starting from a level plane, rise while angling the camera down so the object stays centered, until you are directly above it. Repeat for every side and edge.
  \item Flip the object upside down and repeat the whole process to fully capture the underside.
\end{enumerate}
\figph{Method~1 capture path for a small/medium object (level pass, rising angled-down pass, flipped pass).}

\textbf{Method 2.}
\begin{enumerate}[leftmargin=1.6em]
  \item Walk around the object, camera level, a photo every 5--10\textdegree{} for the full 360\textdegree.
  \item Raise the camera about 15\textdegree{} and angle it down; capture the top following the same 5--10\textdegree{} spacing for a full 360\textdegree.
  \item Raise the camera another 15\textdegree{}, angle it further down, and take another 360\textdegree{} pass (this batch captures more of the top surface).
  \item Flip the object and repeat to fully capture the bottom, so the two sets can be combined.
\end{enumerate}

\noindent\textit{Rule of thumb:} if every other rule is followed, you can combine both methods to capture even more detailed photos.

\textbf{Backgrounds:} Best results occur with limited objects or colors in the surroundings; it doesn't have to be empty, as long as the main object is the focus and meets the 70--80\% fill rule.

\textbf{Shine:} Apply the shine-reduction guidance from Section~\ref{sec:common-shine} as needed.

\subsection{RealityScan}
\inout{organized photos of the object}{an initial 3D model (textured and colorized)}
\begin{todobox}{To be documented}
Note any RealityScan steps or settings specific to small/medium objects that differ from the common workflow (Section~\ref{sec:common-rs}).
\end{todobox}

\subsection{Rigging (Blender)}
\inout{an initial 3D model of the textured, colorized object}{a finalized, textured, decimated, rigged object ready for Unity}
\begin{todobox}{To be documented}
Note any rigging steps specific to small/medium objects that differ from the common workflow (Section~\ref{sec:common-rig}).
\end{todobox}

% ============================================================
\section{Large Objects}
\textit{Example: a desk.} Follow the Common Workflow (Section~\ref{sec:common}); only the differences are listed below. The main difference is that large objects often can't be flipped, so an alternative bottom-capture pass is provided.

\subsection{Photography}
\inout{camera and lighting set up per Section~\ref{sec:common-photo}}{organized photos of the object, split into Top and Bottom folders}

\textbf{Focal length:} 25--40\,mm, with around 35\,mm used most often and giving good results.

\paragraph{Capturing the object.} As before, keep the object locked in place and fully inside the frame.

\textbf{Method 1 (if possible).}
\begin{enumerate}[leftmargin=1.6em]
  \item Walk around the object, camera level, a photo every 5--10\textdegree{} for the full 360\textdegree.
  \item Rise from a level plane while angling the camera down to keep the object centered, until directly above it; repeat for every side and edge.
  \item If the object can be flipped, flip it and repeat the full process for the underside.
  \item \textbf{If it can't be flipped} (common for large objects): again start from a level plane, but this time \emph{lower} while angling the camera \emph{up} to keep the object centered, repeating for every side and edge to capture the bottom.
\end{enumerate}
\figph{Method~1 capture path for a large object, including the no-flip ``lower and angle up'' bottom pass.}

\textbf{Method 2.}
\begin{enumerate}[leftmargin=1.6em]
  \item Walk around the object, camera level, a photo every 5--10\textdegree{} for the full 360\textdegree.
  \item Raise the camera about 10--15\textdegree{}, angle down, and take a full 360\textdegree{} pass of the top.
  \item Raise another 10--15\textdegree{}, angle further down, and take another 360\textdegree{} pass (more top surface).
  \item If the object can be flipped, flip it and repeat to capture the bottom.
  \item \textbf{If it can't be flipped:} after finishing the top, lower the camera about 10--15\textdegree{} and angle it up, taking a full 360\textdegree{} pass of the bottom; then lower another 10--15\textdegree{}, angle further up, and take another 360\textdegree{} pass (more bottom surface).
\end{enumerate}

\textbf{Backgrounds:} Same as Section~\ref{sec:common-photo}: limited objects/colors, with the object as the focus meeting the 70--80\% fill rule.

\textbf{Shine:} Apply the shine-reduction guidance from Section~\ref{sec:common-shine} as needed.

\subsection{RealityScan}
\inout{organized photos of the object}{an initial 3D model (textured and colorized)}
\begin{todobox}{To be documented}
Note any RealityScan steps or settings specific to large objects that differ from the common workflow (Section~\ref{sec:common-rs}).
\end{todobox}

\subsection{Rigging (Blender)}
\inout{an initial 3D model of the textured, colorized object}{a finalized, textured, decimated, rigged object ready for Unity}
\begin{todobox}{To be documented}
Note any rigging steps specific to large objects that differ from the common workflow (Section~\ref{sec:common-rig}).
\end{todobox}

% ============================================================
\section{Extra Large Objects (Room)}
Follow the Common Workflow (Section~\ref{sec:common}); the differences below are significant, because RealityScan is meant for specific objects rather than whole rooms, so room models will be limited and act mainly as a foundation.

\subsection{Photography}
\inout{camera and lighting set up for an enclosed space}{organized photos of the room}

\textbf{Lighting:} Use flat, even lighting from the ceiling as much as possible. Do not let light come in from doors or windows.

\textbf{Focal length:} 20--35\,mm, with the 30--35\,mm range being best to reduce distortion.

\textbf{Focus:} Don't focus on any single point in the room; keep the photos ``flat'' so that as many details as possible for every object in the room are captured equally.

\textbf{Filling the frame:} Capture as much of the room as possible at each level.

\textbf{Orientation:} Shoot vertical to maximize the information in each photo. It takes more time than horizontal, but horizontal limits how much of the top and bottom of the room you see.

\paragraph{Capturing the room.} Make sure windows aren't visible (shut curtains or cover them) and that nothing is moving, with all loose objects secured.
\begin{enumerate}[leftmargin=1.6em]
  \item Stand in the middle of the room (or the area you're capturing). Take a photo, turn slightly left by about 5--10\textdegree, and take another. Repeat for a full 360\textdegree.
  \item Angle the camera up to capture the upper walls and roof, again a photo every 5--10\textdegree{} for a full 360\textdegree.
  \item Angle the camera down to capture the lower walls and floor, again a photo every 5--10\textdegree{} for a full 360\textdegree.
\end{enumerate}
At this point a room model can be rendered to act as a starting point for mimicking rooms.
\figph{Three-pass room capture from a central standing position (level, angled up, angled down).}

\paragraph{Adding detail (optional).} For more detail, take photos above/in front of the furniture, circling the room with 60--80\% overlap per photo so RealityScan knows where each photo is angled relative to the others. If you want something specific captured (e.g., the inside of an open bag), photograph everything in and around it, with transitions between the ``major'' angles. Don't be afraid to capture from slightly different angles to show how an object looks on all sides.

\paragraph{Reducing glare in a room.} Capturing rooms doesn't allow the normal use of scanning spray, dry shampoo, or cross-polarization, so reduce glare in this order:
\begin{enumerate}[leftmargin=1.6em]
  \item Remove as many reflective objects as possible.
  \item Cover any reflective objects that can't be moved (e.g., drape a blanket over a fixed mirror).
  \item As a last, labor-intensive option, coat remaining reflective objects with scanning spray.
  \item For anything that still can't be moved or covered, limit the number of photos that include it.
\end{enumerate}
If there are too many mirrored objects you can't deal with, the room likely can't be captured with 2D photos alone, and LiDAR scanning (to sense object depths) becomes the better option.

\subsection{RealityScan}
\inout{organized photos of the room}{an initial 3D model (textured and colorized)}
\begin{todobox}{To be documented}
Note any RealityScan steps or settings specific to rooms that differ from the common workflow (Section~\ref{sec:common-rs}).
\end{todobox}

\subsection{Rigging (Blender)}
\inout{an initial 3D model of the textured, colorized room}{a finalized, textured, decimated, rigged room ready for Unity}
\begin{todobox}{To be documented}
Note any rigging steps specific to rooms that differ from the common workflow (Section~\ref{sec:common-rig}).
\end{todobox}

% ============================================================
\section{Person}
Follow the Common Workflow (Section~\ref{sec:common}); the key difference is that RealityScan really only works for non-moving subjects, so the entire challenge is keeping the person still.

\subsection{Photography}
\inout{camera and lighting set up around a seated person}{organized photos of the person}

\textbf{Focal length:} 35--50\,mm, with around 35\,mm used most often and giving good results. Don't zoom; keep the person in focus.

\textbf{Filling the frame:} Capture the chest and above, filling about 70--80\% of the frame. Position yourself based on your focal length to meet this.

\textbf{Orientation:} Horizontal or vertical depending on the situation.

\paragraph{Capturing the person.} The person must stay completely still, blinking between photos at most, for the full duration (about 5--10 minutes). Seat them in a flat chair with arms resting on their legs to limit movement.
\begin{enumerate}[leftmargin=1.6em]
  \item Start at the person's face, camera level, taking photos every 5--10\textdegree{} (following the usual photo guidelines).
  \item Raise the camera about 10--15\textdegree{} for a higher angle, again every 5--10\textdegree.
  \item Raise another 10--15\textdegree{} for an even higher angle, again every 5--10\textdegree.
  \item Return to the level plane, then lower the camera about 10--15\textdegree{} for a lower angle; this pass only needs to happen once.
\end{enumerate}
Work as efficiently as possible to limit movement. If the person clearly moves (repositions an arm, itches, smiles, etc.), the process has to restart so that as many characteristics as possible stay the same.
\figph{Capture angles for a seated person (face-level pass, two raised passes, one lowered pass).}

\textbf{Jewelry / shine:} If the person is wearing shiny jewelry, try to have them remove it. If that's undesired, use the shine- and glare-reduction methods from Section~\ref{sec:common-shine} (scanning spray or cross-polarization).

\subsection{RealityScan}
\inout{organized photos of the person}{an initial 3D model (textured and colorized)}
\begin{todobox}{To be documented}
Note any RealityScan steps or settings specific to people that differ from the common workflow (Section~\ref{sec:common-rs}).
\end{todobox}

\subsection{Rigging (Blender)}
\inout{an initial 3D model of the textured, colorized person}{a finalized, textured, decimated, rigged model ready for Unity}
\begin{todobox}{To be documented}
Note any rigging steps specific to people that differ from the common workflow (Section~\ref{sec:common-rig}).
\end{todobox}

% ============================================================
\section{Importing into Unity}
\inout{a finalized, textured, decimated, rigged model from any category above}{the model in use within Unity for VR/simulation applications}
\begin{todobox}{To be documented, final shared step}
This is the final step and is the same for every type of scan, so it lives here as a single shared section. Document the Unity import workflow once: importing the model and its textures/materials, scale and orientation setup, applying the rig, and any settings needed for VR use.
\end{todobox}

% ============================================================
\section{Future Research}
\label{sec:future}

\subsection{Scanning Spray and Cross-Polarization}
KCNSC has approved scanning spray, but we couldn't access it and had to substitute dry shampoo, which reduced glare at the cost of color accuracy. From research, the most professional way to remove glare is cross-polarization: place polarizing sheets over the light sources and a circular polarizing filter on the lens, rotated 90\textdegree{} relative to the lights. This removes the glare.

\subsection{Turntable}
\begin{todobox}{To be expanded}
Investigate using a turntable to standardize the capture rotation. (No notes recorded yet.)
\end{todobox}

\subsection{Lighting}
The lighting requirements were clear, but much of the testing used less-than-ideal lighting and still succeeded. Further research into exactly how much great lighting improves the photos, to find the sweet spot between convenience and fully capturing the object, would be ideal.

\subsection{LiDAR with RealityScan}
RealityScan supports LiDAR, along with photos to support it. Combining LiDAR and photos to build a model could be a major step toward more accurate models and may also reduce the glare problem: the LiDAR data creates a point cloud (like a model) while the photos support it.

\subsection{Photography Setups (Including Object Elevation)}
The best setups would have great lighting covered by polarizing films and multiple cameras aimed at the object (potentially 100+ cameras to capture everything at once), using measured angles so each camera aligns perfectly and the object is captured from many angles at a single moment in time. This would reduce errors from framing (every camera can be set up and checked for perfect framing) and from movement (RealityScan hates movement, but capturing everything in one moment means there is no measurable movement). It would especially excel at capturing humans, whose micro-movements are nearly invisible to the eye but are significant for RealityScan models.

% ============================================================
\section{Conclusion}
\begin{todobox}{Draft, review and refine}
Starting draft assembled from the project notes; expand once the RealityScan, rigging, and Unity sections are written.
\end{todobox}

\noindent This project produced a practical, repeatable pipeline for creating Unity-ready 3D models from ordinary 2D photographs, without relying on AI or cloud services. By organizing the workflow around object scale (small/medium objects, large objects, rooms, and people) and by following a consistent set of photography rules (sufficient overlap, matte surfaces, controlled lighting, and consistent subjects), accessible equipment can produce models suitable for VR and simulation use. The largest remaining limitations involve reflective surfaces and capturing whole rooms, both of which point toward LiDAR and improved capture setups as natural next steps for future work.

\end{document}